% =====================================================================
%  CHAPTER 4 — DESIGN OF THE INTERFEROMETRY SIMULATION PIPELINE
% =====================================================================
\chapter{DESIGN OF THE INTERFEROMETRY SIMULATION PIPELINE}
\label{ch:pipeline}

Chapter~\ref{ch:theory} established the theoretical chain that connects solar
radio flux to interferometric imaging: flux density is converted to received
power, received power to voltage, voltage to correlated visibility, and
visibility to a dirty image. Chapter~\ref{ch:antenna} then provided the
receiving-element model required by this chain, including the antenna gain,
primary beam, input impedance, and mutual-coupling validation. This chapter
implements those two foundations as a reproducible, end-to-end simulation
pipeline, and it is organised so that every subsection can be read as the
direct implementation of one specific relation from Chapter~\ref{ch:theory}.

The purpose of the pipeline is not to reproduce the internal hardware of the
Learmonth observatory. Learmonth provides the external solar input
$S_\nu(t)$ at \SI{610}{\mega\hertz} (Section~\ref{subsec:sv-of-t}); the
simulated instrument is the proposed low-cost four-element UHF
interferometer characterised in Chapter~\ref{ch:antenna}. Each stage in this
chapter therefore answers a specific forward-modelling question: given the
solar flux incident on Earth, what signal would be received by the modelled
Yagi element, what voltage would appear in each receiver channel, what
visibility would be produced by each antenna pair, and what dirty image
would be reconstructed from the resulting $uv$-coverage?

% =====================================================================
\section{From Theory to Implementation: Simulation Chain}
\label{sec:blockdiagram}

The theoretical chain of Chapter~\ref{ch:theory} can be written as one
sequence of quantities:
\begin{equation*}
  S_\nu(t) \;\rightarrow\; P_\mathrm{sig}(t) \;\rightarrow\; V_\mathrm{rms}(t)
  \;\rightarrow\; v_k(t) \;\rightarrow\; V_{ij}(u,v,w) \;\rightarrow\; I_D(l,m).
\end{equation*}
This chapter is the implementation of exactly that sequence, expanded to
include the two steps-real archived data and a real antenna beam-that
Chapter~\ref{ch:theory} treated abstractly:
\begin{equation*}
  \text{Learmonth flux file} \rightarrow \text{simulation time grid}
  \rightarrow \text{primary beam} \rightarrow \text{RF power}
  \rightarrow \text{complex voltage} \rightarrow \text{noise}
  \rightarrow \text{geometric phase} \rightarrow \text{correlation}
  \rightarrow \text{visibility} \rightarrow \text{dirty image}.
\end{equation*}
Table~\ref{tab:theory-to-pipeline} makes this correspondence explicit,
equation by equation, and each subsection below opens with a sentence
identifying which relation it implements.

\begin{table}[H]
  \centering
  \caption[Correspondence between the Chapter~2 theoretical chain and its Chapter~4 implementation]{Correspondence between the Chapter~\ref{ch:theory} theoretical chain and its Chapter~\ref{ch:pipeline} implementation.}
  \label{tab:theory-to-pipeline}
  \renewcommand{\arraystretch}{1.3}
  \begin{tabularx}{\textwidth}{l >{\raggedright\arraybackslash}X >{\raggedright\arraybackslash}X}
    \toprule
    \textbf{Chapter~\ref{ch:theory} gives} & \textbf{Physical meaning} & \textbf{Chapter~\ref{ch:pipeline} implements as} \\
    \midrule
    $S_\nu(t)$, in sfu (Section~\ref{subsec:sv-of-t}) & Solar flux density & Learmonth \SI{610}{\mega\hertz} input file (Section~\ref{sec:data}) \\
    $1~\sfu=\SI{e-22}{\watt\per\square\metre\per\hertz}$ (Eq.~\eqref{eq:sfu}) & Conversion to SI & \texttt{flux\_sfu * 1e-22} (Section~\ref{sec:data}) \\
    $A_e=\lambda^2 G/4\pi$ (Eq.~\eqref{eq:aeff}) & Effective aperture & Computed from the 4NEC2 boresight gain (Section~\ref{sec:beam-integration}) \\
    $A_n(\theta,\phi)$ (Section~\ref{subsec:antenna-gain-beam}) & Primary beam & Loaded from the 4NEC2 pattern of Chapter~\ref{ch:antenna} (Section~\ref{sec:beam-integration}) \\
    $P_\mathrm{sig}=\tfrac12 A_e S_\nu \Delta\nu$ (Eq.~\eqref{eq:friis}) & Received signal power & Flux-to-power block (Section~\ref{sec:voltage-synth}) \\
    $P_n=k_B T_\mathrm{sys}\Delta\nu$ (Eq.~\eqref{eq:johnson}) & Thermal noise power & Noise injection (Section~\ref{sec:voltage-synth}) \\
    $V_\mathrm{rms}=\sqrt{P_\mathrm{rx}R_\mathrm{load}}$ (Eq.~\eqref{eq:vrms}) & Power-to-voltage & Voltage synthesis (Section~\ref{sec:voltage-synth}) \\
    $\phi_k=2\pi\nu\tau_k$ (Eq.~\eqref{eq:taug}) & Geometric phase & Per-antenna phase model (Section~\ref{sec:geometry}) \\
    $V_{ij}=\langle v_i v_j^*\rangle$ (Eq.~\eqref{eq:correlator-def}) & Visibility & Correlator (Section~\ref{sec:correlator}) \\
    $(u,v,w)$ (Eq.~\eqref{eq:uvw}) & Baseline in spatial-frequency coordinates & Geometry module (Section~\ref{sec:geometry}) \\
    $I_D=I*B$ (Eq.~\eqref{eq:dirty}) & Dirty-image relation & Imaging and PSF validation (Section~\ref{sec:imaging}) \\
    \bottomrule
  \end{tabularx}
\end{table}

The pipeline is implemented in Python as a sequence of physically grounded,
independently testable stages, structured around two complementary
simulation engines that share identical geometry, beam, and noise models but
differ in how they produce visibilities. \emph{Engine~A}
(\texttt{simulate\_voltages} $\rightarrow$ \texttt{correlate}) synthesises
per-antenna complex voltage time series, applies geometric phase
sample-by-sample, and time-averages the cross-products to form visibilities:
it is the more physical of the two, following every link of the chain above
explicitly, including realistic signal-to-noise ratios and the
radiometer-equation improvement with integration time
(Section~\ref{subsec:radiometer}). \emph{Engine~B}
(\texttt{forward\_visibilities}) applies the van~Cittert-Zernike relation of
Eq.~\eqref{eq:vcz} analytically to a parametric sky model (a uniform solar
disk), evaluating the resulting visibility function directly at the tracked
$uvw$ coordinates. Engine~B bypasses the voltage and noise chain entirely,
making it orders of magnitude faster for the array-geometry sweeps of
Section~\ref{sec:cases} where absolute noise levels are not the quantity of
interest. Both engines share the same $\mathrm{ENU}\to XYZ\to uvw$ geometry
transformation (Section~\ref{sec:geometry}) and produce the same visibility
format, so the downstream imaging code of Section~\ref{sec:imaging} is
engine-agnostic.

\begin{figure}[H]
  \centering
  \includegraphics[width=\textwidth]{fig_pipeline_block}
  \caption[Block diagram of the end-to-end simulation pipeline]{Block diagram of the end-to-end simulation pipeline.\\[4pt]
  Real SWS Learmonth solar flux enters Engine~A (voltage-synthesis path), which
  constructs per-antenna complex baseband voltages via the Friis reception
  relation, adds Johnson-Nyquist thermal noise, imposes the geometric phase delay
  through the ENU$\rightarrow$XYZ$\rightarrow\uvw$ transformation, and
  accumulates cross-correlations to form visibilities. Engine~B bypasses voltage
  synthesis and evaluates the van Cittert-Zernike forward integral analytically,
  producing visibilities directly from a sky brightness model at negligible
  computational cost. Both engines deliver identically formatted visibility sets
  to the common imaging stage, which computes the dirty image and the
  synthesised beam by vectorised direct DFT.}
  \label{fig:pipeline-block}
\end{figure}

% =====================================================================
\section{Input Data: Learmonth 610~MHz Flux as the External Solar Reference}
\label{sec:data}

This section implements the flux-density input $S_\nu(t)$ introduced in
Section~\ref{subsec:sv-of-t}. As emphasised there and in
Section~\ref{subsec:learmonth-hardware}, the simulation does not attempt to
reproduce the internal receiver chain of the Learmonth observatory; it uses
the Learmonth record as an external, already-calibrated flux-density input,
and forward-models the voltage, noise, and visibilities of the proposed
instrument from that input.

\subsection{What the Learmonth Data Represents}
\label{subsec:learmonth-represents}

As described in Section~\ref{subsec:learmonth-hardware}, Learmonth Solar
Observatory is part of the RSTN, and its \SI{610}{\mega\hertz} channel is one
of eight discrete-frequency Radio Interference Monitoring Set (RIMS)
channels, fed by an \SI{8.5}{\metre} parabolic
dish~\citep{sws,sws_rims_metadata}, digitised at \SI{1}{\hertz} and reported
as calibrated total power in solar flux units~\citep{rstn_overview}. Each
daily file stores one integrated flux sample per second across the local
daytime window, with nighttime samples represented by a six-character
missing-value token (\texttt{//////}). At the Learmonth latitude
$\varphi=-22.2^\circ$, slightly more than half of the \num{86400} seconds in
a UTC day fall in this nighttime gap; because the observatory operates at
UTC$+8$, the genuine daytime window spans roughly 22:20~UTC of the previous
day to 09:30~UTC of the current day, giving approximately \SI{10.4}{\hour} of
real solar samples per day, with a fill factor close to \SI{99}{\percent}
after nighttime removal.

In the notation of this chapter, the raw record is a set of timestamped
samples $\{t_i^\mathrm{raw}, S_i^\mathrm{raw}\}$, with $S_i^\mathrm{raw}$
given in solar flux units. This is exactly $S_\nu(t)$ of
Section~\ref{subsec:sv-of-t}, sampled at Learmonth's native one-second
cadence rather than on the simulation's own time grid.

\subsection{What the Learmonth Data Does Not Provide}
\label{subsec:learmonth-limits}

To avoid overstating what the Learmonth record can justify, it is worth being
explicit about its limits, consistent with the caveat already raised in
Section~\ref{subsec:learmonth-hardware}:
\begin{itemize}[leftmargin=1.5em,itemsep=1pt,topsep=2pt]
  \item Learmonth provides a \emph{calibrated, standardised total-power} solar
        flux density in sfu. It does \emph{not} provide a raw voltage
        waveform.
  \item Learmonth is a single-dish, total-power radiometer at each frequency,
        not an interferometer; it provides \emph{no} phase information
        between separate antennas.
  \item The public metadata do not specify the complete RF/IF filter
        response, gain chain, detector implementation, or calibration
        procedure of the Learmonth receiver (Section~\ref{subsec:learmonth-hardware}).
        These internal details are therefore not assumed anywhere in this
        thesis.
\end{itemize}
Consequently, this thesis does not simulate the Learmonth receiver. It uses
$S_\nu(t)$ as an external solar flux input and forward-models the voltage,
noise, phase, and visibilities that the \emph{proposed} low-cost
interferometer-not Learmonth's own hardware-would produce in response to
that flux. The remainder of this chapter is the implementation of that
forward model.

\subsection{Preprocessing and Interpolation onto the Simulation Grid}
\label{subsec:preprocessing}

A dedicated preprocessing module, \texttt{sws\_to\_csv.py}, handles every
aspect of raw file ingestion robustly. The module exposes a single public
function, \texttt{load\_day(date, freq, out\_dir)}, whose responsibilities
are:
\begin{enumerate}[leftmargin=*]
  \item \textbf{Timestamp preservation.} The parser reconstructs a proper UTC
  timestamp for every sample, retaining all eight RSTN frequencies in a wide
  DataFrame. A \texttt{gap\_flag} column marks any sample that was missing
  or interpolated, so that downstream code can distinguish genuine solar data
  from filled gaps.

  \item \textbf{Nighttime block removal.} Rather than forward-filling the last
  valid daytime sample across the entire night-which would smear a single
  flux value across thousands of seconds and corrupt the transit statistics-
  the function identifies the contiguous nighttime block from the leading and
  trailing runs of \texttt{//////} tokens and excises them entirely. Only the
  genuine daytime segment is returned.

  \item \textbf{CSV export.} The cleaned daily record is exported to a
  \texttt{.csv} file in the project's \texttt{file\_excel\_du\_lieu/} directory
  so that the raw data can be inspected in a spreadsheet at any time, without
  re-running the pipeline.
\end{enumerate}

The function call \texttt{sws.load\_day(NGAY, freq="F610", out\_dir=CSV\_DIR)}
returns two arrays: \texttt{t\_raw} (a timestamp index) and \texttt{sfu\_raw}
(the flux density in Solar Flux Units), covering only the daytime transit
window. Figure~\ref{fig:sws-flux} shows the result for the validation day,
20~March~2025.

\begin{figure}[H]
  \centering
  \includegraphics[width=\textwidth]{fig_sws_flux_day}
  \caption[SWS Learmonth 610~MHz solar flux record for 20~March~2025]{SWS Learmonth 610~MHz solar flux record for 20~March~2025.\\[4pt]
  After preprocessing by \texttt{sws\_to\_csv.py}: nighttime \texttt{//////}
  gap tokens removed, gap flags applied, and only the genuine daytime transit
  window retained. The quiet-Sun continuum is visible as a slowly varying
  envelope; impulsive spikes are candidate Type~III or Type~I burst events.
  This one-sample-per-second record is subsequently interpolated onto the
  50\,000-point simulation time grid before voltage synthesis.}
  \label{fig:sws-flux}
\end{figure}

The raw SWS record contains one sample per second, whereas the simulation
voltage chain operates on a uniformly spaced grid of $t_n=t_0+n\Delta t$,
$n=0,\ldots,N-1$, $N=\num{50000}$, spanning the same transit window. A
\texttt{scipy} cubic spline (\texttt{CubicSpline}) resamples the daytime flux
record onto this grid:
\begin{equation}
  S_\nu(t_n) = \mathrm{interp}\bigl(\{t_i^\mathrm{raw}\},
               \{S_i^\mathrm{raw}\},\,t_n\bigr), \quad n = 0,\ldots,N-1,
  \label{eq:spline}
\end{equation}
where $t_i^\mathrm{raw}$ and $S_i^\mathrm{raw}$ are the original Learmonth
timestamps and flux samples, $t_n$ is the simulation time-grid sample, and
any negative interpolated value is clipped to zero (physically, flux is
non-negative). \emph{This interpolation does not introduce new solar
physics.} It is only a resampling step that places the external solar flux
record onto the same time grid used by the voltage, geometry, and
correlation modules below; the cubic spline is used in preference to linear
interpolation only because it preserves burst morphology and local slope
better and avoids the Runge oscillations of a global polynomial fit.

% =====================================================================
\section{Antenna Response Imported from Chapter~\ref{ch:antenna}}
\label{sec:beam-integration}

This section implements the primary beam $A_n(\theta,\phi)$ of
Section~\ref{subsec:antenna-gain-beam}, replacing the idealised Gaussian
primary beam that is standard in textbook treatments with the \emph{actual}
radiation pattern of the HKD~G11KD-T2 Yagi antenna extracted from the 4NEC2
antenna simulator in Chapter~\ref{ch:antenna}
(Section~\ref{sec:primary-beam-extraction}). The motivation is that the
primary beam directly enters the signal amplitude (Eq.~\eqref{eq:psig}) and
thus affects every downstream result: a Gaussian approximation can
underestimate off-axis attenuation by several dB and overestimate the usable
transit window.

\subsection{Primary-Beam File and Coordinate Convention}
\label{subsec:beam-class}

The primary beam is encapsulated in a \texttt{PrimaryBeam} Python class that
loads the pre-computed \texttt{UFRAME\_primary\_beam\_610MHz.npz} file
produced at the end of Chapter~\ref{ch:antenna}
(Section~\ref{sec:primary-beam-extraction}) by the post-processing script
\texttt{parse\_nec\_out.py}. The file stores three arrays: \texttt{theta\_deg}
(polar angle grid), \texttt{phi\_deg} (azimuth grid), and \texttt{gain\_dBi}
(the two-dimensional gain surface in dBi at \SI{610}{\mega\hertz}, on the
$181\times361$, $1^\circ$-step grid described in
Section~\ref{sec:primary-beam-extraction}). Upon construction the class:
\begin{enumerate}[leftmargin=*]
  \item Locates the boresight automatically as the grid point of maximum gain
  (no assumed direction is hard-coded), using \texttt{np.nanargmax}.
  \item Computes the normalised power response
  $A_n(\theta,\phi) = 10^{[G(\theta,\phi)-G_{\rm max}]/10}$, so that the
  boresight response is exactly unity and the pattern is expressed as a
  fractional power multiplier-the same $A_n$ used in the visibility integral
  of Eq.~\eqref{eq:vcz}.
  \item Constructs a \texttt{RegularGridInterpolator} over the
  $(\theta,\phi)$ grid (bilinear, extrapolating to zero outside the defined
  domain) for fast evaluation at arbitrary angles.
\end{enumerate}

\subsection{Beam Weighting Along the Solar Track}
\label{subsec:beam-weighting}

Two public methods are provided:
\begin{itemize}
  \item \texttt{response\_power($\theta$, $\phi$)}: returns the normalised
  power response $A_n(\theta,\phi)$ at an arbitrary direction.
  \item \texttt{response\_voltage\_offset($\Delta\theta$)}: returns the
  \emph{voltage amplitude} response $\sqrt{A_n}$ as a function of the angular
  offset from boresight in the boresight azimuth plane. This is the method
  called during voltage synthesis (Section~\ref{sec:voltage-synth}), where
  the Sun's angular distance $\theta_z(t_n)$ from the beam peak at each time
  step is computed from the observation geometry of
  Section~\ref{sec:geometry}.
\end{itemize}
The class reports the extracted peak gain $G_{\rm max}$ in dBi on
initialisation; for the HKD~G11KD-T2 this is the value derived from the
4NEC2 antenna simulator (\SI{9.5}{dBi}) and is used to compute the effective
aperture $A_e = G_{\rm max}\,\lambda^2/(4\pi)$ of Eq.~\eqref{eq:aeff} in the
Friis chain of Section~\ref{sec:voltage-synth}.

% =====================================================================
\section{Engine~A: Flux-to-Voltage Forward Model}
\label{sec:voltage-synth}

This section implements Sections~\ref{sec:flux-to-voltage}
and~\ref{sec:noise} of Chapter~\ref{ch:theory}: it carries the interpolated
flux density $S_\nu(t_n)$ of Eq.~\eqref{eq:spline} through the Friis
reception relation and the voltage-equivalent relation
(Eqs.~\eqref{eq:friis}-\eqref{eq:vrms}) to a complex per-antenna voltage with
correctly modelled thermal noise (Eqs.~\eqref{eq:johnson}-\eqref{eq:noise-voltage}).
Table~\ref{tab:params} collects, in one place, every symbol used in this
section and the next, so that no quantity appears in an equation before it
has been defined.

\begin{table}[H]
  \centering
  \caption[Symbols used in the flux-to-voltage and correlation stages]{Symbols used in the flux-to-voltage and correlation stages (Sections~\ref{sec:voltage-synth}-\ref{sec:correlator}).}
  \label{tab:params}
  \renewcommand{\arraystretch}{1.25}
  \begin{tabularx}{\textwidth}{l >{\raggedright\arraybackslash}X l}
    \toprule
    \textbf{Symbol} & \textbf{Meaning} & \textbf{Unit} \\
    \midrule
    $t_i^\mathrm{raw}$, $S_i^\mathrm{raw}$ & Original Learmonth sample time and flux & s (UTC), sfu \\
    $t_n$ & Simulation time-grid sample & s (UTC) \\
    $S_\nu(t_n)$ & Interpolated flux density, SI units & \si{\watt\per\square\metre\per\hertz} \\
    $f_c$ & Centre (observing) frequency & Hz \\
    $\lambda=c/f_c$ & Wavelength & m \\
    $\Delta\nu$ & Simulated receiver bandwidth & Hz \\
    $G_\mathrm{max}$ & Peak antenna gain (Chapter~\ref{ch:antenna}) & dBi / linear \\
    $A_e$ & Effective aperture & m$^2$ \\
    $A_n(\theta,\phi)$ & Normalised primary beam (Chapter~\ref{ch:antenna}) & dimensionless \\
    $\theta_z(t_n)$ & Solar offset from boresight at $t_n$ & deg \\
    $R_\mathrm{load}$ & Receiver load resistance & $\Omega$ \\
    $T_\mathrm{sys}$ & Assumed system temperature of the \emph{simulated} receiver & K \\
    $P_\mathrm{sig}(t_n)$ & Received solar signal power & W \\
    $P_n$ & Thermal noise power & W \\
    $V_\mathrm{rms}(t_n)$ & RMS signal voltage & V \\
    $\sigma_n$ & Noise-voltage standard deviation (per I/Q component) & V \\
    $v_k(t_n)$ & Complex baseband voltage, antenna $k$ & V \\
    $\bm{r}_k$ & ENU position vector of antenna $k$ & m \\
    $\hat{\bm{s}}(t_n)$ & Unit vector toward the Sun & dimensionless \\
    $\tau_k(t_n)$ & Geometric delay, antenna $k$ & s \\
    $\phi_k(t_n)$ & Geometric phase, antenna $k$ & rad \\
    $V_{ij}$ & Complex visibility, baseline $(i,j)$ & complex \\
    $u,v,w$ & Baseline coordinates & wavelengths \\
    \bottomrule
  \end{tabularx}
\end{table}

\subsection{System Parameters and Units}
\label{subsec:sys-params}

All physical constants and instrument parameters are declared in a single
configuration cell~(Cell~0) and inherited by every downstream stage without
being re-declared. Table~\ref{tab:sim-params} summarises the full parameter
set; note that $T_\mathrm{sys}=\SI{150}{\kelvin}$ is an assumed effective
system temperature for the \emph{simulated} receiver being designed in this
thesis-it is not, and is not intended to be, the system temperature of the
Learmonth receiver (Section~\ref{subsec:learmonth-limits}).

\begin{table}[H]
  \centering
  \caption[Global simulation parameters (notebook Cell~0)]{Global simulation parameters (notebook Cell~0).\\[4pt]
  Every physical constant and instrument parameter used throughout the pipeline
  is declared once in this cell and inherited by all downstream stages without
  repetition. Changing a value here (observation date, system temperature,
  bandwidth, or array geometry) propagates automatically to every figure, table,
  and exported file without modifying any other cell.}
  \label{tab:sim-params}
  \renewcommand{\arraystretch}{1.25}
  \begin{tabularx}{\textwidth}{l l >{\raggedright\arraybackslash}X}
    \toprule
    \textbf{Parameter} & \textbf{Symbol / Variable} & \textbf{Value} \\
    \midrule
    Centre frequency        & \texttt{fc}               & \SI{610}{\mega\hertz} \\
    Instantaneous bandwidth & \texttt{BW\_Hz}           & \SI{1}{\mega\hertz} \\
    Load resistance         & \texttt{R\_load}          & \SI{75}{\ohm} \\
    System temperature      & \texttt{T\_sys}           & \SI{150}{\kelvin} \\
    Simulation samples      & \texttt{N\_SAMPLES}       & \num{50000} \\
    Correlator dump size    & \texttt{INTEGRATION\_BINS}& 50 \\
    Site latitude           & \texttt{LAT\_DEG}         & $-22.2^\circ$ (Learmonth~\citep{learmonth_coords}) \\
    Solar declination       & \texttt{DEC\_SUN\_DEG}    & $0.0^\circ$ (20 Mar 2025, spring equinox~\citep{jpl_horizons}) \\
    Transit window          & \texttt{H\_deg}           & $[-90.0^\circ,\,+90.0^\circ]$ \\
    Clock offsets           & \texttt{CLOCK\_OFFSET\_NS} & $[0,0,0,0]$~ns (zero: coherent baseline; swept in Sec.~\ref{sec:clock-offset}) \\
    \bottomrule
  \end{tabularx}
\end{table}

\subsection{Signal Power from Flux Density}
\label{subsec:signal-power}

This subsection implements the Friis reception relation of
Eq.~\eqref{eq:friis}. The flux density $S_\nu(t_n)$ interpolated at
Eq.~\eqref{eq:spline} is first converted from Solar Flux Units to SI
(1~SFU $= 10^{-22}$~W\,m$^{-2}$\,Hz$^{-1}$, Eq.~\eqref{eq:sfu}) and then to a
reference received signal power, evaluated at the peak antenna gain:
\begin{equation}
  P_{\mathrm{sig}}(t) = \tfrac{1}{2}\,A_e\,S_\nu(t)\,\Delta\nu,
  \label{eq:psig}
\end{equation}
where $A_e$ is the peak effective aperture computed from the 4NEC2 boresight
gain ($A_e = G_{\rm max}\lambda^2/4\pi$, Eq.~\eqref{eq:aeff}), and
$\Delta\nu = \SI{1}{\mega\hertz}$ is the assumed instantaneous receiver
bandwidth of the simulated system (Table~\ref{tab:params}). $P_\mathrm{sig}$
is not yet a voltage, and it does not yet include the antenna's directional
response away from boresight; both are added next.

\subsection{Equivalent Voltage Synthesis}
\label{subsec:voltage-synth-detail}

This subsection implements the voltage-equivalent relation of
Eq.~\eqref{eq:vrms}. The root-mean-square signal voltage at peak gain is
\begin{equation}
  V_{\mathrm{rms}}(t) = \sqrt{P_{\mathrm{sig}}(t)\,R_{\mathrm{load}}},
\end{equation}
directly from $P_\mathrm{sig}(t)$ of Eq.~\eqref{eq:psig}. The antenna's
directional response then enters at this same voltage level, through the
primary beam of Section~\ref{sec:beam-integration}: the beam-weighted signal
envelope used in the complex voltage model of Eq.~\eqref{eq:vk} below is
\begin{equation}
  A_\mathrm{env}(t_n) = V_\mathrm{rms}(t_n)\cdot\sqrt{A_n(\theta_z(t_n))} .
\end{equation}
It is worth being explicit about what this quantity is and is not. It is the
voltage that the \emph{simulated} Yagi element and receiver, described in
Chapter~\ref{ch:antenna}, would produce if the incident solar flux were
$S_\nu(t)$. It is \emph{not} a voltage recovered from Learmonth hardware:
Learmonth reports only a calibrated flux density
(Section~\ref{subsec:learmonth-limits}), never a receiver voltage, so there
is nothing to invert or recover. This voltage is generated forward, from
flux to power to voltage, exactly as Section~\ref{sec:flux-to-voltage}
prescribes.

\subsection{Receiver Noise Model}
\label{subsec:noise-model}

This subsection implements the Johnson-Nyquist noise relation of
Eq.~\eqref{eq:johnson} and the noise-voltage relation of
Eq.~\eqref{eq:noise-voltage}. Even though the flux density $S_\nu(t)$
produces a measurable signal power, no real channel observes that signal
alone: every channel also carries thermal noise from the antenna, the sky,
the ground, the cable, and the receiver electronics. The simulation must
therefore generate not only the deterministic solar voltage of
Section~\ref{subsec:voltage-synth-detail}, but also a statistically correct
noise voltage to add to it.

Each antenna receives an independent additive Johnson-Nyquist noise
contribution~\citep{nyquist1928} of power $P_n = k_B\,T_{\rm sys}\,\Delta\nu$
(Eq.~\eqref{eq:johnson}), where, as noted in
Section~\ref{subsec:sys-params}, $T_\mathrm{sys}$ is an assumed effective
system temperature for the simulated receiver, not a value measured from
Learmonth. In complex baseband, this noise power is represented as a
circularly symmetric Gaussian process: the noise standard deviation per I/Q
component is
\begin{equation}
  \sigma_n = \frac{\sqrt{P_n\,R_{\mathrm{load}}}}{\sqrt{2}},
  \label{eq:noise-sigma}
\end{equation}
matching Eq.~\eqref{eq:noise-voltage}, where the $\sqrt{2}$ distributes the
total noise power equally between the in-phase and quadrature components.
The independence of thermal noise across antennas is enforced by using a
fresh draw from \texttt{np.random.default\_rng} for each antenna inside
\texttt{simulate\_voltages}.

\subsection{Complex Voltage and Geometric Phase}
\label{subsec:complex-voltage}

This subsection implements the correlator input voltage of
Section~\ref{subsec:visibility-amp-phase} together with the geometric delay
and phase of Eq.~\eqref{eq:taug}. Before the equation is written down, every
symbol in it is defined: $k$ indexes the antenna; $\bm{r}_k$ is the ENU
position vector of antenna $k$ (Section~\ref{sec:geometry}); $\hat{\bm{s}}(t_n)$
is the unit vector toward the Sun at simulation time $t_n$; the geometric
delay of antenna $k$ is $\tau_k(t_n)=\bm{r}_k\cdot\hat{\bm{s}}(t_n)/c$
(Eq.~\eqref{eq:taug}); and the corresponding geometric phase is
$\phi_k(t_n)=2\pi f_c\,\tau_k(t_n)$.

A coherent far-field source process $\varepsilon(t)$ is drawn once from a
unit-variance complex Gaussian distribution with a fixed random seed (42), so
that all simulation runs with the same configuration are reproducible and
antenna responses can be compared on a sample-by-sample basis. The per-antenna
complex baseband voltage is then
\begin{equation}
  v_k(t_n) = \underbrace{A_{\mathrm{env}}(t_n)\,\varepsilon(t_n)}_{%
    \text{signal}} \,e^{\,i\phi_k(t_n)}
    \;+\; \underbrace{n_k(t_n)}_{\text{noise}},
  \label{eq:vk}
\end{equation}
where $A_\mathrm{env}(t_n)$ is the beam-weighted signal amplitude of
Section~\ref{subsec:voltage-synth-detail}, $\phi_k(t_n)$ is the geometric
phase of antenna~$k$ defined above and derived in full in
Section~\ref{sec:geometry}, and $n_k(t_n)$ is the per-antenna noise sample of
standard deviation $\sigma_n$ from Eq.~\eqref{eq:noise-sigma}. The result is
collected into an $(N_{\mathrm{samp}}\times N_{\rm ant})$ complex matrix
\texttt{V\_matrix}.

\subsection{Clock Offset and Quantisation}
\label{subsec:clock-offset}

This subsection covers two remaining engineering steps that act on the
voltage of Eq.~\eqref{eq:vk}: a controllable clock-offset perturbation used
to study phase-coherence loss (Chapter~\ref{ch:intro},
Section~\ref{subsec:element-vs-system}), and the hardware-export
quantisation path.

\paragraph{Clock offset.} In a hardware interferometer each digitiser is
driven by its own oscillator. Without a common phase reference, each antenna
accumulates an independent timing offset $\tau_k$ (in seconds) relative to a
notional master. At the carrier frequency $f_c = \SI{610}{\mega\hertz}$ this
offset is indistinguishable from a geometric delay and manifests in complex
baseband as a constant phase rotation applied to every voltage sample of
that antenna:
\begin{equation}
  v_k(t_n)\;\longrightarrow\;
  v_k(t_n)\,\exp\!\bigl(j2\pi f_c\,\tau_k\bigr),
  \label{eq:clock-phase}
\end{equation}
where $\tau_k$ is expressed in seconds so that the exponent $2\pi f_c\tau_k$ is
dimensionless. The simulation parameter \texttt{CLOCK\_OFFSET\_NS} accepts the
same quantity in nanoseconds and scales by $10^{-9}$ internally before applying
Eq.~\eqref{eq:clock-phase} to each column of \texttt{V\_matrix}.
The default value is $[0,0,0,0]$~ns (Table~\ref{tab:sim-params}), recovering
the perfectly phase-coherent baseline. For baseline $(i,j)$, the differential
offset $\tau_{ij}=\tau_j-\tau_i$ shifts the measured visibility by an extra
phase without affecting its amplitude:
\begin{equation}
  \hat{V}_{ij,\,\mathrm{meas}}
  = \hat{V}_{ij,\,\mathrm{true}}\,
    \exp\!\bigl(j2\pi f_c\,\tau_{ij}\bigr).
  \label{eq:vis-clock}
\end{equation}
A static offset is equivalent to an uncalibrated phase error in the
visibility, removable by a point-source calibration step but corrupting if
left uncorrected. When $\tau_{ij}$ \emph{drifts} during an integration dump
because the two oscillators run at slightly different rates, the
time-averaged fringe amplitude is reduced, a process known as
\emph{fringe washing}-precisely the effect diagnosed in the Xuan~Mai field
campaign of Section~\ref{subsec:element-vs-system}. The quantitative analysis
is given in Section~\ref{sec:clock-offset} of Chapter~\ref{ch:results}.

\paragraph{Signal-to-noise validation.} At the end of the voltage synthesis
stage a sanity check computes the peak signal-to-noise ratio at solar transit
($\theta_z = 0$ for a transit directly overhead) in decibels:
\begin{equation}
  \mathrm{SNR}_{\mathrm{peak}} = 10\log_{10}\!\left(
    \frac{\max_t P_{\mathrm{sig}}(t)\cdot R_{\mathrm{load}}}{P_n\cdot R_{\mathrm{load}}}
  \right).
  \label{eq:snr-peak}
\end{equation}
This is printed alongside $A_e$, $T_{\rm sys}$, and $\Delta\nu$ to confirm
that the signal sits above the noise floor by a physically reasonable margin
before any correlation gain is applied. Because \texttt{V\_matrix} contains
\num{50000} complex samples, plotting the raw $|v(t)|$ waveform directly is
uninformative: a 500-sample sliding-RMS filter is therefore applied,
\begin{equation}
  \hat{V}_{\mathrm{rms}}[n] = \sqrt{\frac{1}{W}\sum_{k=n}^{n+W-1}|v[k]|^2},
  \quad W=500,
  \label{eq:sliding-rms}
\end{equation}
computed efficiently via a cumulative sum (O$(N)$). Figure~\ref{fig:voltage-env}
shows the resulting envelope: a correctly synthesised \texttt{V\_matrix} should
show a smooth bell-shaped envelope peaking at mid-transit, with the noise
floor $V_\mathrm{noise}=\sqrt{P_n R_\mathrm{load}}$ visible in the wings.

\begin{figure}[H]
  \centering
  \includegraphics[width=\textwidth]{fig_voltage_env}
  \caption[Per-antenna voltage envelope (sliding RMS) over the transit window]{Per-antenna voltage envelope (sliding RMS) over the transit window.\\[4pt]
  The 500-sample sliding root-mean-square of the complex voltage at one
  representative antenna (all four are statistically identical under this
  beam and geometry model) traces the combined effect of the solar flux
  profile and the primary-beam attenuation as the Sun rises, transits, and
  sets. The dashed line marks the Johnson-Nyquist noise floor
  $V_{\rm noise}\approx0.394\,\si{\micro\volt}$. Per-sample SNR is sub-unity
  (the signal is buried in noise), but the bell-shaped transit envelope is
  unambiguous on the scale of the sliding RMS and serves as the first
  pipeline diagnostic checkpoint.}
  \label{fig:voltage-env}
\end{figure}

\paragraph{16-bit quantisation and hardware export.} Although an FPGA
hardware correlator was removed from the project scope, the voltage-chain
engine retains a full 16-bit quantisation and hex-export stage (Cell~4 of the
notebook), both to demonstrate forward compatibility with real-time hardware
and to provide an additional numerical validation of the voltage amplitudes:
a misconfigured signal chain would either saturate the ADC or bury the
signal in quantisation noise. The full-scale range of a 16-bit signed
integer is $[-32768, 32767]$, so a global scale factor is computed as
\begin{equation}
  k_s = \frac{32767}{\max_{n,k}|v_k(t_n)|},
  \label{eq:scale}
\end{equation}
mapping the peak voltage across all antennas and time steps to the ADC
ceiling. Each complex sample is then rounded to the nearest integer, clamped,
and represented as a two's complement 16-bit word for both the in-phase (I)
and quadrature (Q) components. Before writing the output file, a validation
check confirms that the noise floor occupies at least 3 least-significant
bits:
\begin{equation}
  \mathrm{LSB}_{\mathrm{noise}} = V_{\mathrm{noise}}\cdot k_s \geq 3.
  \label{eq:lsb-check}
\end{equation}
If this condition fails, the thermal noise would be below the quantisation
floor and the digital correlator would see quantisation noise rather than
Johnson-Nyquist noise, corrupting the noise statistics; for the parameters of
Table~\ref{tab:sim-params} the condition is comfortably satisfied. The
quantised output is written to a plain-text file
\texttt{iq\_array\_stimulus.txt}, with one row per time sample and
$2N_{\rm ant}$ hexadecimal columns per row (I and Q for each antenna, in
order), together with a self-documenting metadata file
\texttt{iq\_array\_stimulus\_meta.txt} recording every parameter needed to
decode the file.

% =====================================================================
\section{Array Geometry and \texorpdfstring{$uvw$}{uvw} Sampling}
\label{sec:geometry}

This section implements the geometry of Section~\ref{sec:uvw-coverage}: the
ENU ground coordinates, the source-based $uvw$ transformation of
Eq.~\eqref{eq:uvw}, and the Earth-rotation $uv$-coverage argument of
Section~\ref{subsec:earth-rotation-uv}.

\subsection{Array Configuration Presets}
\label{subsec:array-presets}

Antenna positions are specified in the local East-North-Up (ENU) frame of
Section~\ref{subsec:enu}, with Up $= 0$ for all antennas (a flat
ground-based deployment); this is the $\bm{r}_k$ of
Section~\ref{subsec:complex-voltage}. The pipeline provides six configurable
array constructors, all accessed through a common \texttt{Array} class
interface:

\begin{table}[H]
  \centering
  \caption[Array geometry presets in the simulation pipeline (Cell~B)]{Array geometry presets in the simulation pipeline (Cell~B).\\[4pt]
  All six constructors share the same ENU$\rightarrow$XYZ$\rightarrow\uvw$
  geometry engine and produce output in the identical format consumed by the
  correlator and imaging stages. The \texttt{from\_enu(coords)} constructor
  accepts an arbitrary list of antenna positions, enabling any custom layout
  to be evaluated without modifying the pipeline code.}
  \label{tab:array-presets}
  \renewcommand{\arraystretch}{1.25}
  \begin{tabularx}{\textwidth}{l >{\raggedright\arraybackslash}X l}
    \toprule
    \textbf{Constructor} & \textbf{Description} & \textbf{Baseline range} \\
    \midrule
    \texttt{preset\_compact()}         & 4 E-W antennas at 0, 8, 16, 24~m          & 8--24~m \\
    \texttt{preset\_extended()}        & Graded 2-D cross, 5--130~m baselines       & 5--130~m \\
    \texttt{linear(8, 15, az=0)}       & 8 antennas due North, 15~m spacing         & 15--105~m \\
    \texttt{tee(6, 5, 18)}             & 6 E-W + 5 N antennas, 18~m spacing        & 18--162~m \\
    \texttt{random\_2d(16,120,seed=3)} & 16 antennas uniformly random, $120\times120$~m & variable \\
    \texttt{from\_enu(coords)}         & Arbitrary list of (E, N) coordinates       & arbitrary \\
    \bottomrule
  \end{tabularx}
\end{table}

For Case~1 (the compact field array), \texttt{preset\_compact()} places four
antennas along the East-West axis at positions
$(E,N) = (0,0),\,(8,0),\,(16,0),\,(24,0)$~m, yielding baselines of 8, 16, and
24~m ($N_B = 3$ pairs). The \texttt{Array} object automatically computes all
baseline vectors, pair indices, and baseline lengths, which are used by the
correlator and imaging stages without any further user input.

\subsection{Geographic Observation Context}
\label{subsec:geo-context}

The observation window is parameterised by the solar hour angle
$H \in [-90.0^\circ,\,+90.0^\circ]$, which corresponds to the interval
over which the Sun is above the local horizon at Learmonth
($\varphi = -22.2^\circ$~\citep{learmonth_coords}) for 20~March~2025
(spring equinox; solar declination $\delta = 0.0^\circ$~\citep{jpl_horizons}). At the equinox the Sun rises and sets exactly on the
equator, so the observable window is exactly $\SI{12.0}{\hour}$, consistent
with the SWS transit record for that date.

Each hour-angle step corresponds to $\approx 278$ simulation
samples per degree of Earth rotation. Figure~\ref{fig:globe-track} shows the
Sun's apparent track rendered on an orthographic globe centred over the
Australia--Southeast~Asia region; the orange band marks the full
$\SI{12.0}{\hour}$ observing window ($H\in[-90^\circ,+90^\circ]$) and the
yellow star marks the transit instant ($H=0$).

\begin{figure}[H]
  \centering
  \includegraphics[width=\textwidth]{fig_globe_track}
  \caption[Apparent solar track and transit geometry (Learmonth, 20~March~2025)]{Apparent solar track and transit geometry over the Learmonth--Southeast~Asia
  region on 20~March~2025 ($\delta=0.0^\circ$, spring equinox).\\[4pt]
  The orange band marks the 12-hour observing window ($H\in[-90^\circ,+90^\circ]$,
  sunrise to sunset). The transit point (yellow star, $H=0$) is located on the
  terrestrial equator at longitude $114.1^\circ$~E, midway between
  Hanoi ($21.0^\circ$~N, $105.9^\circ$~E) and Learmonth ($22.2^\circ$~S,
  $114.1^\circ$~E) in geographic latitude.
  %
  The transit point is determined by two independent facts: (1)~transit ($H=0$)
  is defined as the instant the Sun crosses the \emph{observer's} meridian
  ($\lambda_\text{obs} = 114.1^\circ$~E), fixing the longitude; and
  (2)~the sub-solar latitude equals the solar declination $\delta=0^\circ$,
  placing the sub-solar point on the equator. Hanoi lies
  $8.2^\circ$ west of the Learmonth meridian; the Sun therefore passed through
  the Hanoi longitude $\approx 33$~minutes \emph{before} culminating at
  Learmonth, which is why the transit star does not appear directly above
  Hanoi.}
  \label{fig:globe-track}
\end{figure}

\subsection{The ENU \texorpdfstring{$\rightarrow$}{->} XYZ \texorpdfstring{$\rightarrow$}{->} \texorpdfstring{$\uvw$}{(u,v,w)} Transformation}
\label{subsec:enu2xyz}

The first rotation converts each antenna's local East-North-Up displacement
$(\Delta E,\,\Delta N,\,\Delta U)$-the ENU baseline of
Section~\ref{subsec:enu}-into the equatorial baseline frame
$(X_\lambda,Y_\lambda,Z_\lambda)$ used by Eq.~\eqref{eq:uvw}.
At observing latitude $\varphi$ this rotation is \citep{tms2017}
\begin{equation}
  \begin{bmatrix} X_\lambda \\ Y_\lambda \\ Z_\lambda \end{bmatrix}
  = \frac{1}{\lambda}
  \begin{bmatrix}
    -\sin\varphi & 0 & \cos\varphi \\
     0           & 1 & 0           \\
     \cos\varphi & 0 & \sin\varphi
  \end{bmatrix}
  \begin{bmatrix} \Delta E \\ \Delta N \\ \Delta U \end{bmatrix},
  \label{eq:enu2xyz}
\end{equation}
where the factor $1/\lambda$ converts the physical baseline (metres) to
dimensionless wavelengths. For the flat deployment considered here
$\Delta U = 0$, so the $\cos\varphi$ column is inactive.

Antenna positions are rotated first from ENU to the equatorial XYZ frame by
the site latitude (Eq.~\eqref{eq:enu2xyz}), and then, for each hour angle, into
the baseline coordinate frame by Eq.~\eqref{eq:uvw}. A single call to
\texttt{ix.compute\_uvw(baseline\_vectors\_enu, obs)} returns the
$(N_I\times N_B)$ arrays $u$, $v$, $w$ for all integration dumps and all baselines
simultaneously, where $N_I$ is the number of correlator dump times. The same
transformation supplies all three quantities consistently: the per-antenna
geometric phase $\phi_k = -\frac{2\pi}{\lambda}\bm{r}_k\cdot\hat{\bm{s}}$
used in Eq.~\eqref{eq:vk}, the baseline $\uvw$ coordinates for imaging, and
the $w$-term for fringe stopping.

A validation step cross-checks the dump nearest to transit
($H \approx 0^\circ$) against an independent single-point call
\texttt{compute\_uvw($H$=0)}, confirming agreement in $u$ to within 1~$\lambda$
and in $v$ to within 0.5~$\lambda$ across all baselines.

\subsection{\texorpdfstring{$uv$}{uv} Coverage}
\label{subsec:uv-coverage}

As $H$ sweeps across the transit, Earth rotation causes each baseline vector to
trace an arc in the $uv$-plane, exactly as described in
Section~\ref{subsec:earth-rotation-uv}. For the compact east-west array the
arcs are short and nearly horizontal (small $v$ component), yielding a
one-dimensional sampling line in the $uv$-plane; for two-dimensional layouts
the arcs curve and fill the plane more completely. Figure~\ref{fig:uv-coverage}
shows the $uv$-coverage of the compact array.

\begin{figure}[H]
  \centering
  \includegraphics[width=\textwidth]{fig_uv_coverage}
  \caption[$uv$-coverage of the compact east-west array over the 12-hour transit]{$uv$-coverage of the compact east-west array over the 12-hour transit.\\[4pt]
  Case~1, 20~March~2025, $H\in[\pm90.0^\circ]$). Each of the three baselines
  (8~m, 16~m, 24~m) traces a nearly horizontal arc as the Earth rotates.
  Together the three arcs sample only a narrow east-west strip of the
  $uv$-plane, providing angular resolution only in the east-west direction.
  Recovering north-south resolution requires either additional north-south
  baselines or observation at non-equatorial declinations.}
  \label{fig:uv-coverage}
\end{figure}

% =====================================================================
\section{Correlator and Visibility Formation}
\label{sec:correlator}

This section implements the correlator definition of
Section~\ref{subsec:visibility-amp-phase}, Eq.~\eqref{eq:correlator-def}. A
single antenna voltage $v_k(t_n)$, as synthesised in
Section~\ref{sec:voltage-synth}, is not yet an interferometric measurement:
a single antenna can only measure power, $\langle|v_k(t)|^2\rangle$, which
carries no information about the direction of the source on the sky. The
interferometric measurement is the time-averaged cross-product between two
antenna voltages.

\subsection{Cross-Correlation and Time Averaging}
\label{subsec:cross-corr}

For each antenna pair $(i,j)$ the correlator accumulates the time-averaged
product of the complex voltage streams:
\begin{equation}
  \hat{V}_{ij}[m] = \frac{1}{N_{\mathrm{dump}}}
    \sum_{n\in\mathcal{D}_m} v_i(t_n)\,v_j^{*}(t_n),
  \label{eq:correlate}
\end{equation}
where:
\begin{itemize}[leftmargin=1.5em,itemsep=1pt,topsep=2pt]
  \item $i,j$ index the antenna pair;
  \item $\mathcal{D}_m$ is the set of $N_{\mathrm{dump}} = $
        \texttt{INTEGRATION\_BINS} $= 50$ samples in the $m$-th integration
        dump;
  \item $v_j^*$ is the complex conjugate of antenna $j$'s voltage;
  \item $\hat{V}_{ij}[m]$ is the resulting complex visibility estimate for
        dump $m$, with amplitude $|\hat{V}_{ij}|$ and phase
        $\angle\hat{V}_{ij}$.
\end{itemize}
$N_I = N_{\mathrm{samp}}/N_{\mathrm{dump}} = 1000$ is the total number of
visibility dumps. Equation~\eqref{eq:correlate} is the discrete estimator of
the true visibility $V(u_{ij},v_{ij})$ at the instantaneous baseline of the
pair, exactly as defined in Eq.~\eqref{eq:correlator-def}. The integration
time per dump is
\begin{equation}
  \tau_{\mathrm{dump}} = N_{\mathrm{dump}}\cdot\frac{12\times 3600\,\mathrm{s}}{N_{\mathrm{samp}}}
                  \approx \SI{43.2}{\second},
\end{equation}
short enough that time-average smearing across the \SI{3}{\degree} FOV is
negligible at the antenna separations used here
(Section~\ref{subsec:earth-rotation-uv}).

\subsection{Correlator Output: Amplitude and Phase}
\label{subsec:corr-output}

The correlator returns an $(N_I \times N_B)$ complex array of visibilities.
Post-processing extracts the amplitude $|{\hat V}_{ij}|$ and phase
$\angle\hat{V}_{ij}$ for all baselines simultaneously. A 10-dump sliding-mean
filter is applied to the amplitude to reduce the sample-to-sample noise scatter
while preserving the transit bell curve:
\begin{equation}
  \overline{|V_{ij}|}[m] = \frac{1}{10}\sum_{k=m}^{m+9}|\hat{V}_{ij}[k]|.
\end{equation}
The van~Cittert-Zernike theorem (Eq.~\eqref{eq:vcz}) predicts that, for an
incoherent extended source, longer baselines measure lower correlated power
because they resolve more structure in the brightness distribution: at the
transit peak, $|\hat{V}_{ij}|_{\rm transit}$ must be a monotonically
non-increasing, jinc-shaped function of baseline length
(Eq.~\eqref{eq:jinc}). The pipeline prints a per-baseline amplitude table at
the transit bin and flags any baseline that violates the expected ordering
with a \texttt{WARN} label. This check would fail if, for example, the
geometric phase were applied with the wrong sign or the source process were
not fully coherent across antennas, and it provides confidence that the
signal chain from Eq.~\eqref{eq:vk} through Eq.~\eqref{eq:correlate} is
physically correct. Case~1's four short, closely-spaced baselines are too
narrow a range to trace out the jinc envelope, so this validation is shown
empirically once wider baselines are available, alongside the extended
Case~2 array of Section~\ref{subsec:case2} (Fig.~\ref{fig:visibility-jinc}).

% =====================================================================
\section{Imaging: Dirty Image and PSF}
\label{sec:imaging}

This section implements the imaging relations of
Section~\ref{sec:visibility-to-image}: the direct Fourier inversion of
Eq.~\eqref{eq:dirty-ft} and the convolution relation of Eq.~\eqref{eq:dirty}.

\subsection{Vectorised Direct Discrete Fourier Transform}
\label{subsec:dft-imaging}

The dirty image is formed by the exact inversion of the van~Cittert-Zernike
integral (Eq.~\eqref{eq:dirty}) on a uniform $(l,m)$ sky grid:
\begin{equation}
  I^D(l_p, m_q) = \sum_{b=1}^{N_B}\sum_{t=1}^{N_I}
    \hat{V}_b[t]\,e^{+2\pi i(u_b[t]\,l_p + v_b[t]\,m_q)},
  \label{eq:dirty-dft}
\end{equation}
where $(l_p, m_q)$ are direction cosines on the image grid, and
$u_b[t], v_b[t]$ are the tracked $uv$-coordinates of baseline~$b$ at dump~$t$.
The sum is evaluated exactly for every visibility, avoiding any gridding
approximation. The implementation uses a vectorised NumPy \texttt{einsum}
over baselines and pixels to exploit BLAS-level parallelism; for the number of
visibilities used here ($\lesssim 10^4$) this direct approach is both accurate
and fast enough.

Fringe stopping (removal of the $w$-term) is applied when
\texttt{do\_fringe\_stop=True}: each visibility is phase-rotated by
$e^{-2\pi i\,w_b[m](1-\sqrt{1-l^2-m^2})}$ before the DFT sum. Fringe
stopping is enabled for all Engine~A outputs (where the voltages carry the full
$w$ phase) and disabled for Engine~B outputs (where the forward model already
evaluates on the fringe-stopped baseline).

\subsection{The Synthesised (Dirty) Beam}
\label{subsec:dirty-beam-impl}

The synthesised beam (PSF) is computed from the identical sampling function by
replacing all visibilities with unity before Eq.~\eqref{eq:dirty-dft}:
\begin{equation}
  B^D(l,m) = \sum_{b,t} e^{+2\pi i(u_b\,l + v_b\,m)}.
  \label{eq:psf-dft}
\end{equation}
The PSF is normalised so that its peak equals unity. Confirming that
$I^D = I^{\rm sky} \ast B^D$ (Eq.~\eqref{eq:dirty}) by comparing the dirty
image with the input sky convolved independently with the dirty beam is the
principal internal validation of the imaging stage. Figure~\ref{fig:sky-conv-psf}
shows this verification for the compact array.

\begin{figure}[H]
  \centering
  \includegraphics[width=\textwidth]{fig_sky_conv_psf}
  \caption[Dirty image verification: DFT inversion versus sky convolved with PSF]{Dirty image verification: DFT inversion versus sky convolved with PSF.\\[4pt]
  Left: the true point-source sky model, prior to any convolution. Centre:
  that sky model convolved with the independently computed synthesised beam.
  Right: the dirty image produced by the direct DFT inversion of the
  simulated visibilities (Eq.~\eqref{eq:dirty-dft}). Pixel-level numerical
  agreement between the centre and right panels would confirm that the
  correlator, the $uv$-sampling function, and the inverse DFT are mutually
  consistent and together implement Eq.~\eqref{eq:dirty} correctly.}
  \label{fig:sky-conv-psf}
\end{figure}

\begin{quote}\small\itshape
Note: in the current rendering, the centre panel (Sky~$\ast$~PSF) shows a
smooth, sidelobe-free brightness gradient, while the right panel (dirty
image from the DFT) shows a clear sidelobe pattern. These are not yet
pixel-for-pixel consistent and should be re-verified against the notebook
before this figure is treated as a completed validation.
\end{quote}

\subsection{FWHM Measurement}
\label{subsec:fwhm}

The angular resolution is characterised by the full-width at half-maximum of
the dirty beam measured separately in the East-West ($l$) and North-South ($m$)
directions. One-dimensional cuts through the PSF centre are extracted:
\begin{equation}
  B^{\rm EW}(l) = B^D(l,\,m\!=\!0), \quad
  B^{\rm NS}(m) = B^D(l\!=\!0,\,m).
\end{equation}
The FWHM in each direction is computed as the width of the above-half-maximum
region: $\Theta_{\rm FWHM} = l_{\rm right} - l_{\rm left}$ where
$B \geq B_{\rm max}/2$. These values are printed and annotated on the PSF
panel of every evaluation case. The ratio $\Theta_{\rm NS}/\Theta_{\rm EW}$
is used as a summary metric of coverage anisotropy.

\begin{figure}[H]
  \centering
  \begin{subfigure}{0.48\textwidth}
    \includegraphics[width=\textwidth]{fig_dirty_beam}
    \caption{Synthesised dirty beam (PSF)}
  \end{subfigure}\hfill
  \begin{subfigure}{0.48\textwidth}
    \includegraphics[width=\textwidth]{fig_dirty_image}
    \caption{Dirty image of the Sun}
  \end{subfigure}
  \caption[Dirty beam and dirty image for the compact east-west array (Case~1)]{Dirty beam and dirty image for the compact east-west array (Case~1).\\[4pt]
  Left: the synthesised dirty beam (point-spread function) is highly elongated
  in the north-south direction, the characteristic signature of a purely
  east-west baseline set that provides no north-south spatial-frequency
  sampling. Right: the resulting dirty image of the Sun inherits this
  elongation, appearing as a north-south stripe rather than a resolved disk.
  These artefacts require either additional north-south baselines or
  deconvolution to remove.}
  \label{fig:dirty}
\end{figure}

\begin{quote}\small\itshape
Note: the numeric FWHM annotation printed directly on the current
\texttt{fig\_dirty\_beam} panel has not yet been cross-checked against the
$\Theta_{\rm EW}\approx35'$ analytic estimate above and should be re-verified
against the notebook before this figure is treated as final.
\end{quote}

% =====================================================================
\section{Evaluation Cases for Chapter~\ref{ch:results}}
\label{sec:cases}

The pipeline is exercised across five evaluation cases that together
demonstrate the capability envelope of the simulation and provide the
quantitative results analysed in Chapter~\ref{ch:results}. Cases~1, 1.5,
and~1.75 use Engine~A and real SWS flux; Cases~2 and the sweep use Engine~B with an
analytic solar disk model, which makes it practical to evaluate many
configurations without the memory cost of per-antenna time series.

% - - - - - - - - - - - - - - - - - - - - - - - - - - - - - - - - - -
\subsection{Case 1: Compact East-West Array (Field Configuration)}
\label{subsec:case1}

Case~1 represents the physical deployment configuration: four antennas on an
east-west baseline at 0, 8, 16, and 24~m, matching the field array at
Xuan~Mai. The array is constructed by \texttt{preset\_compact()}, yielding
$N_B = 3$ baselines (Ant~0-1: 8~m; Ant~0-2: 16~m; Ant~0-3: 24~m).

Engine~A runs the full voltage chain on \num{50000} samples, producing a
$(50000\times 4)$ complex voltage matrix. The correlator accumulates 50 samples
per dump, yielding 1000 dump times. The UV coverage (Fig.~\ref{fig:uv-coverage})
shows three short arcs along the $u$-axis, with negligible $v$ extent because
the baselines are purely east-west. The result has two immediate consequences
for imaging:

\begin{enumerate}[leftmargin=*]
  \item \textbf{East-West resolution:} The 24~m maximum baseline at
  \SI{610}{\mega\hertz} ($\lambda = 49.1$~cm) corresponds to $u_{\rm max}
  \approx 48.9\,\lambda$, giving an E-W FWHM of approximately
  $\Theta_{\rm EW} \approx 1/(2u_{\rm max}) \approx 35'$.

  \item \textbf{North-South resolution:} With $v_{\rm max} \approx 0$, there
  is no North-South fringe information. The N-S PSF FWHM spans the entire FOV
  (\SI{180}{\arcminute} for a \SI{3}{\degree} field), and the dirty image is
  correspondingly smeared: any point source would appear as a stripe running the
  full N-S extent of the image.
\end{enumerate}

The N-S-to-E-W FWHM ratio provides the quantitative measure of this asymmetry.
At the spring-equinox declination used throughout this chapter ($\delta=0^\circ$),
$v\equiv0$ exactly for a purely east-west ground baseline at every hour angle
(Fig.~\ref{fig:uv-coverage} shows this directly: the three tracks are flat lines
at $v=0$), so the ratio is not merely large but formally unbounded within the
simulated field of view, consistent with the ``$\gg1$'' entry of
Table~\ref{tab:case-comparison} below. A non-zero $v$
component, and therefore a finite N-S-to-E-W ratio, would only appear for an
east-west baseline observed away from $\delta=0^\circ$. The dirty image and
PSF panels confirm the expected elongation.

% - - - - - - - - - - - - - - - - - - - - - - - - - - - - - - - - - -
\subsection{Case 1.5: Four-Antenna Cross Layout (2-D Baseline Coverage)}
\label{subsec:case15}

Case~1.5 is a direct controlled comparison with Case~1: the same four antennas,
the same spacing $s$, and the same real SWS data (same \texttt{amp\_env},
\texttt{noise\_sigma}, and \texttt{source\_process}), but rearranged into a
cross layout with two antennas along the East-West axis and two along the
North-South axis:
\begin{equation}
  \text{Ant~0}: (-s,0),\quad \text{Ant~1}: (+s,0),\quad
  \text{Ant~2}: (0,-s),\quad \text{Ant~3}: (0,+s),
\end{equation}
where $s$ is the minimum spacing of the Case~1 compact array (8~m). This
produces $N_B = 6$ baselines: two E-W baselines ($2s$), two N-S baselines
($2s$), and two diagonal baselines ($2\sqrt{2}\,s$). The cross layout is
constructed by \texttt{ix.from\_enu([[-s,0],[+s,0],[0,-s],[0,+s]])}.

Because the same \texttt{source\_process} is used for both cases, the signal
component of every antenna voltage is driven by the same coherent realisation
of the solar signal; the only difference is the geometric phase
$\phi_k(t)$, which changes with antenna position. The comparison is therefore
fair: any difference in the resulting dirty image or PSF is attributable
solely to the changed baseline coverage, not to a different noise draw or
signal realisation.

Table~\ref{tab:case-comparison} and Fig.~\ref{fig:dirty-cross} show the PSF and
1-D cuts for Case~1.5. The two N-S baselines add a non-zero $v$ component to
the $uv$ coverage, recovering North-South angular resolution at the cost of
reducing the maximum E-W baseline from $3s$ (Case~1) to $2s$ (Case~1.5). The
FWHM in the N-S direction improves by the ratio $\Theta_{\rm NS}^{(1)}/
\Theta_{\rm NS}^{(1.5)}$, demonstrating that two-dimensional baseline coverage
is necessary and sufficient to recover N-S resolution even with only four
antennas.

\begin{table}[H]
  \centering
  \caption[PSF FWHM: east-west array (Case~1) versus cross layout (Case~1.5)]{PSF FWHM: east-west array (Case~1) versus cross layout (Case~1.5).\\[4pt]
  Both cases use the same four antennas and the same 8~m element spacing; the
  only difference is that two antennas are rotated to point north in Case~1.5.
  The ratio $\Theta_{\rm NS}/\Theta_{\rm EW}$ directly quantifies coverage
  anisotropy: a value near unity indicates equal angular resolution in both
  axes, which is the goal of a two-dimensional interferometer.}
  \label{tab:case-comparison}
  \renewcommand{\arraystretch}{1.25}
  \begin{tabular}{l c c c}
    \toprule
    \textbf{Layout} & \textbf{FWHM E-W} & \textbf{FWHM N-S} &
    \textbf{$\Theta_{\rm NS}/\Theta_{\rm EW}$} \\
    \midrule
    Case~1 (E-W 4-ant, compact)  & $\sim 35'$   & $\gg 100'$ & $\gg 1$ \\
    Case~1.5 (Cross 4-ant, $s=8$~m)  & $\sim 53'$   & $\sim 53'$  & $\approx 1$ \\
    \bottomrule
  \end{tabular}
\end{table}

\begin{figure}[H]
  \centering
  \includegraphics[width=\textwidth]{fig_dirty_from_dft}
  \caption[Dirty image and PSF for the cross layout (Case~1.5)]{Dirty image and PSF for the cross layout (Case~1.5).\\[4pt]
  The four panels show the dirty image (top-left), the synthesised beam
  (top-right), the east-west 1-D PSF cut (bottom-left), and the north-south
  1-D PSF cut (bottom-right). Adding two north-south baselines recovers
  north-south angular resolution and produces a nearly symmetric PSF, at the
  cost of reducing the maximum east-west baseline from $3s$ to $2s$ and thus
  slightly degrading east-west resolution relative to Case~1.}
  \label{fig:dirty-cross}
\end{figure}

\begin{quote}\small\itshape
Note: the current \texttt{fig\_dirty\_from\_dft} rendering for Case~1.5 prints
an East-West FWHM of ``nan'' and a North-South FWHM that does not yet match
the $\sim53'$, $\Theta_{\rm NS}/\Theta_{\rm EW}\approx1$ values of
Table~\ref{tab:case-comparison}. This should be re-verified against the
notebook before the figure and table are treated as final.
\end{quote}

% - - - - - - - - - - - - - - - - - - - - - - - - - - - - - - - - - -
\subsection{Case 1.75: DLITE-Inspired Triangle-Plus-Centre Layout}
\label{subsec:case175}

Cases~1 and~1.5 share the same base element spacing ($s=8$~m) but differ only in
topology (line versus cross). A third topology at the same spacing scale, motivated
directly by the minimum-redundancy design philosophy of the Deployable Low-Band
Ionosphere and Transient Experiment (DLITE)~\citep{helmboldt2021dlite}, is evaluated
as Case~1.75: three antennas placed at the vertices of an equilateral triangle of
side $8$~m, with the fourth antenna at the triangle's centroid. The circumradius of
an equilateral triangle of side $s$ is $R=s/\sqrt{3}\approx4.62$~m, so the array is
constructed by \texttt{ix.from\_enu} with three centre-to-vertex baselines of
$\approx4.6$~m and three triangle-side baselines of $8.0$~m
(Fig.~\ref{fig:case175-geometry}a). Case~1.75 is run with Engine~A on the same
20~March~2025 SWS flux record used throughout this chapter, so the comparison with
Cases~1 and~1.5 is direct.

Unlike the cross layout of Case~1.5, in which the two East-West baselines are
parallel and redundant with each other (and likewise for the two North-South
baselines and the two diagonals), all six baselines of the triangle-plus-centre
layout point in six different directions. Figure~\ref{fig:case175-geometry}b shows
the resulting $uv$-coverage: six distinct, non-overlapping tracks rather than the
three doubled-up tracks of a symmetric cross. This is precisely the design principle
DLITE uses to extract the maximum number of independent spatial-frequency samples
from a minimal number of elements~\citep{helmboldt2021dlite}, rather than adding
antennas along a fixed pair of orthogonal axes.

\begin{figure}[H]
  \centering
  \includegraphics[width=0.92\textwidth]{fig_case175_geometry}
  \caption[DLITE-inspired triangle-plus-centre array geometry and resulting $uv$-coverage (Case~1.75)]{DLITE-inspired triangle-plus-centre array geometry and resulting $uv$-coverage (Case~1.75).\\[4pt]
  (a) Antenna layout: one element at the centroid and three at the vertices of an
  equilateral triangle of side 8~m, giving three centre-to-vertex baselines of
  4.6~m and three triangle-side baselines of 8.0~m. (b) The corresponding
  $uv$-coverage on 20~March~2025: all six baselines trace out distinct, non-redundant
  tracks, in contrast to the doubled-up tracks produced by the symmetric cross of
  Case~1.5.}
  \label{fig:case175-geometry}
\end{figure}

Figure~\ref{fig:case175-eval} shows the resulting dirty image, PSF, and 1-D profile
cuts. The non-redundant baseline set improves the North-South FWHM from
$180'$ for the purely East-West Case~1 array to approximately $19.8'$, a
$9.1\times$ improvement, achieved with the same four antennas and the same base
spacing unit ($s=8$~m) used to construct Cases~1 and~1.5. This is a better
North-South result than the $\sim53'$ obtained by the symmetric cross of Case~1.5
(Table~\ref{tab:case-comparison}), because six non-redundant $uv$-tracks sample the
spatial-frequency plane more densely than four antennas arranged on two orthogonal
axes. The East-West cut, however, does not fall clearly below half-maximum within
the fixed \SI{3}{\degree} field of view (Fig.~\ref{fig:case175-eval}c), so no
reliable East-West FWHM could be extracted by the automated measurement: with a
maximum baseline of only $8$~m-the same scale as the \emph{shortest} baseline in
Case~1-the triangle-plus-centre layout trades East-West resolution for isotropic,
non-redundant coverage rather than improving resolution outright.

\begin{figure}[H]
  \centering
  \includegraphics[width=\textwidth]{fig_case175_eval}
  \caption[Dirty image, PSF, and 1-D profile cuts for the triangle-plus-centre array (Case~1.75)]{Dirty image, PSF, and 1-D profile cuts for the triangle-plus-centre array (Case~1.75).\\[4pt]
  (a) Dirty image; (b) synthesised beam (PSF), annotated with the measured
  North-South FWHM of $19.8'$; (c) East-West profile through the image centre, for
  which no half-maximum crossing exists within the \SI{3}{\degree} field of view;
  (d) North-South profile through the image centre, showing a symmetric
  half-maximum crossing at $\pm9.9'$. The $9.1\times$ improvement in North-South
  FWHM relative to the purely East-West Case~1 array (\SI{180}{\arcminute})
  confirms that all six non-redundant baselines contribute independent
  spatial-frequency information.}
  \label{fig:case175-eval}
\end{figure}

Case~1.75 is included not as a sixth point in the resolution sweep of
Section~\ref{subsec:sweep}, but as a topology intended to inform a future physical
deployment. This case follows the non-redundant triangle-plus-centre topology
inspired by DLITE, but it is not a full DLITE system simulation and does not use
the original DLITE frequency, antenna, or baseline scale: DLITE operates at
\SIrange{30}{40}{\mega\hertz} with LWA-style crossed-dipole antennas on baselines of
several hundred metres~\citep{helmboldt2021dlite,carson2022dlite}, whereas
Case~1.75 is evaluated at \SI{610}{\mega\hertz} with the same $s=8$~m spacing used
throughout this chapter. What carries over from DLITE is the minimum-redundancy
\emph{topology}-the geometric principle that lets exactly four elements-no
more-recover useful two-dimensional angular resolution-not the specific frequency,
antenna, or baseline scale. This case therefore serves as a topology blueprint for
a possible future four-element array upgrade, should a DLITE-inspired layout be
pursued at UHF.

% - - - - - - - - - - - - - - - - - - - - - - - - - - - - - - - - - -
\subsection{Case 2: Extended Two-Dimensional Array with Forward Visibility Model}
\label{subsec:case2}

Case~2 uses the extended graded-cross preset (\texttt{preset\_extended()}),
which places antennas along both cardinal axes with graded spacing in the range
5-130~m, yielding baselines up to \SI{130}{\metre}. At \SI{610}{\mega\hertz}
this corresponds to $u_{\rm max} \approx 265\,\lambda$, providing an E-W
angular resolution of $\Theta_{\rm EW} \approx 11'$-sufficient to spatially
resolve the solar disk whose angular diameter is $\approx 32'$.

Because the extended array has many more baselines than Case~1, the
per-antenna voltage chain of Engine~A would require several hundred megabytes
of RAM at \num{50000} samples. Engine~B is therefore used instead: the forward
visibility model evaluates the van~Cittert-Zernike integral analytically for a
uniform solar disk of angular diameter $0.53^\circ$ at $N_{\rm img} = 600$
uniformly spaced hour angles across the transit:
\begin{equation}
  V^{\rm fwd}(u,v) = \int\!\!\!\int I_{\rm disk}(l,m)\,
    e^{-2\pi i(ul+vm)}\,\mathrm{d}l\,\mathrm{d}m,
\end{equation}
evaluated at each $(u_{b}[m],\,v_{b}[m])$ point. The primary beam power
weighting $A_n^2(\theta_z(t))$ is applied as a time-varying scalar weight
\texttt{time\_weight} that modulates the effective integration time at each
hour angle.

Because Case~2's graded baselines span a much wider range than Case~1's,
they are also the first point in the pipeline where the jinc-shaped
amplitude decay of Section~\ref{sec:correlator} (Eq.~\eqref{eq:jinc}) can be
traced empirically. Figure~\ref{fig:visibility-jinc} shows the transit-peak
visibility amplitude of every Case~2 baseline as a function of length,
overlaid on the analytic $\mathrm{jinc}(u\,\theta_\odot)$ prediction for a
uniform solar disk of angular diameter approximately 32~arcminutes.

\begin{figure}[H]
  \centering
  \includegraphics[width=\textwidth]{fig_visibility_jinc}
  \caption[Visibility amplitude versus baseline length at transit peak, Case~2]{Visibility amplitude versus baseline length at transit peak, Case~2.\\[4pt]
  Simulated visibility amplitudes (points), from the Case~2 forward-model
  array, are overlaid on the analytic $\mathrm{jinc}(u\,\theta_\odot)$
  prediction (solid curve) for a uniform solar disk of angular diameter
  approximately 32~arcminutes. Agreement across all sampled baselines
  confirms that the correlator and the ENU$\rightarrow\uvw$ geometry code
  correctly implement the van Cittert-Zernike theorem, and that the extended
  array has the sensitivity to detect the partially resolved structure of the
  solar disk.}
  \label{fig:visibility-jinc}
\end{figure}

The resulting dirty image is shown in Fig.~\ref{fig:case2-image}. Because the
largest baselines approach and exceed the Nyquist scale of the solar disk,
the visibility amplitude departs significantly from the flat behaviour seen in
Case~1: the jinc-function envelope of the correlated flux drops below 50\% by
$u\approx 130\,\lambda$ for the $0.53^\circ$ disk. This behaviour is the
interferometric signature of a resolved source and constitutes one of the
central observational results of this simulation.

\begin{figure}[H]
  \centering
  \includegraphics[width=\textwidth]{fig_case2_image}
  \caption[Dirty image from the extended graded-cross array (Case~2, $b_{\rm max}=130$~m)]{Dirty image from the extended graded-cross array (Case~2, $b_{\rm max}=130$~m).\\[4pt]
  The array achieves an east-west angular resolution of approximately
  $\Theta_{\rm EW}\approx11'$, sufficient to begin spatially resolving the
  $32'$ solar disk. The forward visibility model (Engine~B) was used to avoid
  the memory cost of a full per-antenna voltage chain; the analytical
  van Cittert-Zernike evaluation enables a fine 600-point hour-angle grid at
  negligible computational overhead. The visibility amplitude drop beyond
  $u\approx130\,\lambda$ is the interferometric signature of a partially
  resolved extended source.}
  \label{fig:case2-image}
\end{figure}

To manage memory during Case~2, the raw data arrays (\texttt{u2, v2, w2, V2})
are saved to a compressed \texttt{.npz} file and then explicitly deleted from
RAM using \texttt{del} and \texttt{gc.collect()} before plotting. The image
is similarly deleted after rendering. This ensures that even on a
low-RAM workstation, the full pipeline can be executed without kernel restart.

% - - - - - - - - - - - - - - - - - - - - - - - - - - - - - - - - - -
\subsection{Sweep: Five Array Geometries}
\label{subsec:sweep}

The final evaluation is a systematic sweep across five representative array
configurations (Cell~I of the notebook), designed to illustrate how baseline
distribution and maximum baseline length jointly determine image quality. The
five configurations evaluated are:

\begin{enumerate}[leftmargin=*]
  \item \textbf{Compact E-W} (\texttt{preset\_compact()}): the 4-element field
  array, $b_{\rm max} = 24$~m. Low resolution, 1-D coverage.

  \item \textbf{Linear N-S, 8 elements} (\texttt{linear(8,15,az=0)}): eight
  antennas due North with 15~m spacing, $b_{\rm max} = 105$~m. Orthogonal to
  Case~1: good N-S resolution, poor E-W resolution without Earth-rotation.

  \item \textbf{Tee-shaped, 11 elements} (\texttt{tee(6,5,18)}): six E-W
  antennas plus five N-S antennas with 18~m spacing, $b_{\rm max}=162$~m.
  Provides simultaneous coverage in both axes at moderate cost.

  \item \textbf{Extended graded cross} (\texttt{preset\_extended()}): the
  Case~2 array, $b_{\rm max} = 130$~m. Highest resolution; irregular spacing
  reduces sidelobe height.

  \item \textbf{Random 2-D, 16 elements} (\texttt{random\_2d(16,120,seed=3)}):
  sixteen antennas placed uniformly at random within a $120\times120$~m square.
  Demonstrates the alternative sidelobe structure produced by non-redundant,
  non-structured baselines.
\end{enumerate}

For each configuration, the pipeline calls \texttt{ix.forward\_visibilities}
to compute visibilities on a $600$-sample hour-angle grid, forms the dirty image
and PSF on a $192\times 192$ pixel grid over a \SI{3}{\degree} FOV, and saves
three-panel figures (dirty image, PSF, $uv$-coverage) as PNG files. A summary
table records the maximum baseline and the dynamic range metric
$\mathrm{DR} = I^D_{\mathrm{peak}}/\sigma_{I^D}$ (peak image value divided by
image RMS) for each configuration. The five-panel montage is assembled by
\texttt{matplotlib} in a $3\times 5$ grid and saved as \texttt{\_montage.png}.

\begin{figure}[H]
  \centering
  \includegraphics[width=\textwidth]{fig_layout_sweep}
  \caption[Sweep montage: dirty images, PSFs, and $uv$-coverage for five array layouts]{Sweep montage: dirty images, PSFs, and $uv$-coverage for five array layouts.\\[4pt]
  Top row: dirty images of the Sun; middle row: synthesised dirty beams (PSFs);
  bottom row: $uv$-coverage track patterns accumulated over the 12-hour transit.
  Columns from left to right: compact east-west (4 antennas, 24~m max baseline),
  linear north-south (8 antennas, 105~m), tee-shaped (11 antennas, 162~m),
  extended graded cross (130~m), and random two-dimensional (16 antennas,
  120~m). The progressive improvement in PSF symmetry and sidelobe suppression
  with increasing baseline diversity and two-dimensional coverage is directly
  visible across the columns.}
  \label{fig:layout-sweep}
\end{figure}

Table~\ref{tab:sweep-summary} summarises the key metrics. The progression from
compact to extended shows a monotonic improvement in angular resolution (lower
E-W FWHM) and generally increasing dynamic range, while the random layout
achieves competitive sidelobe behaviour despite its lower maximum baseline
because the non-redundant spacing suppresses grating lobes.

\begin{table}[H]
  \centering
  \caption[Array sweep summary: maximum baseline, PSF FWHM, and image dynamic range]{Array sweep summary: maximum baseline, PSF FWHM, and image dynamic range.\\[4pt]
  Maximum baseline lengths, approximate east-west PSF full-width at
  half-maximum (estimated as $\lambda/(2b_{\rm max})$), and image dynamic range
  (peak signal divided by RMS background) are listed for each of the five array
  configurations. The progression from compact to extended demonstrates a
  monotonic improvement in angular resolution; the random layout achieves
  competitive sidelobe suppression due to its non-redundant, non-structured
  baseline distribution.}
  \label{tab:sweep-summary}
  \renewcommand{\arraystretch}{1.25}
  \begin{tabular}{l c c c}
    \toprule
    \textbf{Configuration} & \textbf{$b_{\rm max}$ (m)} &
    \textbf{Est.\ FWHM E-W} & \textbf{DR (peak/RMS)} \\
    \midrule
    Compact E-W (4-ant)      & 24  & $\approx 35'$  & lowest \\
    Linear N-S (8-ant)       & 105 & $\approx 14'$  & moderate \\
    Tee (11-ant)              & 162 & $\approx  9'$  & good \\
    Extended cross            & 130 & $\approx 11'$  & good \\
    Random 2-D (16-ant)       & 120 & $\approx 12'$  & good (low sidelobes) \\
    \bottomrule
  \end{tabular}
\end{table}

% =====================================================================
\section{Chapter Summary}
\label{sec:pipeline-conclusion}

The pipeline has been designed and implemented as a faithful, fully
parameterised forward model of a UHF radio interferometer, built as the
direct implementation of the theoretical chain of Chapter~\ref{ch:theory}
and the receiving-element model of Chapter~\ref{ch:antenna}
(Table~\ref{tab:theory-to-pipeline}). The key design decisions and their
physical motivations are:
\begin{itemize}[leftmargin=*]
  \item \textbf{Real primary beam.} The idealised Gaussian pattern was replaced
  by the actual 4NEC2 radiation pattern loaded via the \texttt{PrimaryBeam}
  class, ensuring that the beam-modulated signal amplitude is physically correct
  at all hour angles.

  \item \textbf{External flux, forward-modelled receiver.} Learmonth supplies
  only the calibrated flux density $S_\nu(t)$; every downstream voltage, noise,
  phase, and visibility is generated by a forward model of the \emph{proposed}
  interferometer, not by reproducing Learmonth's own receiver
  (Section~\ref{subsec:learmonth-limits}).

  \item \textbf{Two simulation engines.} Engine~A (voltage chain) provides
  realistic noise statistics and hardware-compatible IQ output for the field
  array; Engine~B (forward visibility) enables rapid parameter sweeps over many
  array configurations that would be intractable with the full voltage chain.

  \item \textbf{16-bit quantisation path.} The LSB-validation check confirms
  that the noise floor occupies the required minimum number of bits,
  demonstrating that the synthesised signal levels are correctly scaled for
  a real ADC.

  \item \textbf{Consistent single-transformation geometry.} All three geometric
  outputs-per-antenna phase, $\uvw$ baselines, and fringe-stopping $w$-term
  -are derived from a single \texttt{compute\_uvw} call, eliminating
  inconsistency errors that arise when geometry is computed piecemeal.

  \item \textbf{Five evaluation cases.} The compact E-W array (Case~1) and its
  cross variant (Case~1.5) isolate the effect of baseline dimensionality on PSF
  shape. The DLITE-inspired triangle-plus-centre layout (Case~1.75) shows that a
  non-redundant topology recovers North-South resolution from only four
  elements. The extended array with the forward model (Case~2) demonstrates
  resolved-source visibility decay. The five-configuration sweep provides a
  comprehensive survey of the resolution-coverage trade-off space.
\end{itemize}

\begin{quote}\itshape
Learmonth gives the incident solar flux; Chapter~\ref{ch:antenna} gives the
antenna response; this chapter forward-models the voltage, visibility, and
image of the proposed interferometer.
\end{quote}

Together these results establish that the simulation faithfully captures the
physics from incident flux to synthesised image. The quantitative outputs of
Chapter~\ref{ch:results} can therefore be interpreted as predictions of the
behaviour of a real instrument with the corresponding array geometry, not as
numerical artefacts of the implementation. The browser-based visualisation
companion that accompanies this pipeline is not part of this physical model;
it is an accessibility tool, described in
Appendix~\ref{sec:appx-html} together with the rest of the project
repository.
\clearpage
